\section{The common normalizer and the fixed-field envelope}
\label{sec:envelope}

The central group-theoretic feature is stronger than equality of two
permutation characters: after the transport in \eqref{eq:left-transport},
the two embedded subgroups share the same index-two normalizer.  Define
\begin{equation}
 N=N_G(\Hplus)=N_G(\Hthree),
 \label{eq:def-N}
\end{equation}
let $\Hzero$ be the third index-two subgroup of $N$ containing the common
derived subgroup, and put
\begin{equation}
 J=\Hplus\cap\Hthree.
 \label{eq:def-J}
\end{equation}
The one-line generators in \cref{app:group-census} make these definitions
entirely independent of Table-of-Marks numbering.

\begin{proposition}[Exact subgroup lattice]
\label{prop:group-lattice}
The five subgroups have the following invariants.
\begin{center}
\begin{tabular}{@{}lrrrrcc@{}}
\toprule
group & order & $[G:H]$ & SmallGroup & $\ToM$ & $\Core_G(H)$ & $N_G(H)$\\
\midrule
$N$       & $324$ & $160$ & $[324,39]$ & $327$ & $1$ & $N$\\
$\Hplus$  & $162$ & $320$ & $[162,11]$ & $301$ & $1$ & $N$\\
$\Hzero$  & $162$ & $320$ & $[162,10]$ & $302$ & $1$ & $N$\\
$\Hthree$ & $162$ & $320$ & $[162,19]$ & $303$ & $1$ & $N$\\
$J$       & $81$  & $640$ & ---        & $266$ & $1$ & $N$\\
\bottomrule
\end{tabular}
\end{center}
Moreover
\begin{equation}
 [N,N]=J,\qquad N^{\mathrm{ab}}\cong C_2\times C_2,
 \qquad N/J\cong\Vfour,
 \label{eq:N-quotient}
\end{equation}
and for distinct $i,j\in\{+,0,3\}$,
\begin{equation}
 H_i\cap H_j=J,\qquad \langle H_i,H_j\rangle=N.
 \label{eq:meet-join}
\end{equation}
\end{proposition}

\begin{proof}
Generate the subgroups from the one-based image arrays in
\cref{app:durable-arrays}.  Closure gives the displayed orders.  Direct
conjugation by every element of the $51840$-element labelled action gives the
normalizers and cores in the table.  The commutator subgroup of $N$ has order
$81$ and its element set equals the generated subgroup $J$; the quotient has
order four and SmallGroup identifier $[4,2]$, hence is $\Vfour$.  The three
index-two subgroups are the inverse images of the three order-two subgroups
of this quotient.  Their pairwise intersections are therefore the kernel
$J$, and the images of any two distinct ones generate $\Vfour$, proving
\eqref{eq:meet-join}.  The supplementary group certificate repeats these
calculations from the durable arrays in two independent implementations.
\end{proof}

Define the fixed fields
\begin{equation}
 M=K^N,\qquad F_i=K^{H_i}\ (i\in\{+,0,3\}),
 \qquad L=K^J.
 \label{eq:fixed-fields}
\end{equation}

\begin{corollary}[Biquadratic diamond]
\label{cor:field-diamond}
The inclusions and degrees are
\[
 \Q\subset M\subset F_i\subset L\subset K,
 \qquad
 [M:\Q]=160,\ [F_i:\Q]=320,\ [L:\Q]=640.
\]
The extension $L/M$ is Galois with group $\Vfour$, and the three $F_i$ are
exactly its quadratic intermediate fields.
\end{corollary}

\begin{proof}
The degree equalities are $[K^H:\Q]=[G:H]$.  Inclusions reverse subgroup
inclusions, so $J\subset H_i\subset N$ gives $M\subset F_i\subset L$.
Since $J\triangleleft N$, Galois correspondence over $M$ identifies
$\Gal(L/M)$ with $N/J$.  The three groups $H_i/J$ are precisely the three
order-two subgroups of $N/J\cong\Vfour$, proving the last assertion.
\end{proof}

\begin{proposition}[Normal closures and automorphisms]
\label{prop:closures-aut}
Each of $M,\Fplus,\Fzero,\Fthree,L$ has normal closure $K$ over $\Q$, and
\[
 \Aut_{\Q}(M)=1,\qquad
 \Aut_{\Q}(F_i)\cong C_2,\qquad
 \Aut_{\Q}(L)\cong\Vfour.
\]
Furthermore $\Fthree=x(K^{\Hminus})$.
\end{proposition}

\begin{proof}
The core of every defining subgroup is trivial by
\cref{prop:group-lattice}; hence the normal closure of each fixed field inside
$K$ is fixed by the trivial subgroup and equals $K$.  Formula
\eqref{eq:normalizer-aut} and the normalizer column give
\[
 N_G(N)/N=1,\qquad N/H_i\cong C_2,\qquad N/J\cong\Vfour.
\]
Finally, $x$ fixes $\Q$ and sends the $\Hminus$-fixed elements to the
$x\Hminus x^{-1}$-fixed elements.  This proves
$x(K^{\Hminus})=K^{\Hthree}$.  It does not put the original embedded minus
field inside this chosen $L$; only its transported conjugate occurs in the
diamond.
\end{proof}

The lattice is now defined independently of any primitive element.  This
order of proof is deliberate: the carriers in the next section identify
explicit generators for fields already determined by Galois correspondence,
rather than defining fields retroactively from a resolvent degree.
